% ---------- פרק 8: סיכום ----------
% לפי פרק 10 במדריך: ממצאים, מוכנות (Design Readiness) ושלוש פעולות מיידיות אחרי אישור הדוח. עמוד אחד.
\chapter{סיכום}\label{ch:summary}

בארגונים גדולים קליטת ספק חדש נמשכת שבועות \cite{hicx2021}, הסיכון הפיננסי גבוה יותר דווקא בחברות פרטיות, שעליהן יש פחות מידע \cite{rapidratings2026}, וגם ארגון ביטחוני גדול אינו יכול להסתמך רק על מה שהספק מוסר \cite{gao2025}. מאמר הדגל \cite{sadeghi2026} מראה שאפשר לשלב מודל שפה בבחירת ספקים רב-קריטריונית, והמאמר הנלווה \cite{zhao2025} מראה שמודלי שפה מזהים סיכוני ספקים מתוך חדשות, אבל אף אחד מהם לא מפיק המלצה על ספק בודד מתוך מידע ציבורי ולא מודד כמה מהמידע חסר (פרק~\ref{ch:literature}). גם בתיעוד הפומבי של המוצרים המסחריים לא נמצא מוצר שמעריך את איכות הספק, מציג ציון כיסוי מידע ומקשר כל טענה למקור (פרק~\ref{ch:market}). מבין שלוש חלופות נבחרה, במטריצת החלטה משוקללת ובניתוח רגישות, מערכת מרובת סוכנים שמעריכה כל ספק בשני חלקים, פרופיל חוזקות ודגלי סיכון, ומחשבת את שער הפסילה בקוד (פרק~\ref{ch:design}).

\section*{מוכנות לשלב המימוש}
הארכיטקטורה, הנוסחאות, מקורות המידע ותוכנית העבודה מוגדרים, ולכל הסיכונים הקריטיים נקבעו פעולות מנע ותוכניות מגירה (פרקים~\ref{ch:workplan} ו-\ref{ch:risks}). המשקלות הסופיים ומאגר התיקוף נקבעים במשימות הראשונות של תוכנית העבודה (\en{T1.1} ו-\en{T1.3}), ולכן אינם מעכבים את תחילת המימוש.

\section*{פעולות מיידיות}
אחרי אישור הדוח יבוצעו שלוש פעולות (איור~\ref{fig:next}): פגישת התנעה עם המנחה האקדמי והמנחה התעשייתי לאישור תוכנית העבודה, פתיחת מאגר \en{GitHub} לקוד, לנתונים ולפרוטוקולי הפגישות, ותחילת המשימות הראשונות בתוכנית העבודה, מילוי ה-\en{AHP} ובניית מאגר התיקוף.

\begin{figure}[!h]
\centering
\begin{tikzpicture}[
  st/.style={draw, rounded corners, fill=blue!6, align=center, font=\footnotesize, text width=3.1cm, minimum height=10mm},
  arr/.style={-{Stealth[length=2mm]}, thick}]
\node[st] (c) at (0,0) {\foreignlanguage{hebrew}{מילוי ה-\en{AHP} ובניית מאגר התיקוף}};
\node[st] (b) at (3.8,0) {\foreignlanguage{hebrew}{פתיחת מאגר \en{GitHub} לקוד ולנתונים}};
\node[st] (a) at (7.6,0) {\foreignlanguage{hebrew}{פגישת התנעה עם שני המנחים}};
\node[st, fill=green!12] (ok) at (11.4,0) {\foreignlanguage{hebrew}{הדוח אושר}};
\draw[arr] (ok) -- (a); \draw[arr] (a) -- (b); \draw[arr] (b) -- (c);
\end{tikzpicture}
\caption{שלוש הפעולות המיידיות אחרי אישור הדוח}\label{fig:next}
\end{figure}
